\section{A sharp primitive bound with a bounded residual census}
We now assume that $G$ is primitive of degree $n$, distinct from $A_n,S_n$.
If its noncommuting diagram is disconnected, the earlier primitive
commuting-factor argument gives $G=T\times T$ on $|T|=n\geq60$ points
and $d\leq2\log_2n<n/2$. Otherwise every adjacent pair fails to commute.
We explicitly retain the following independent theorem of Conder:
a finite nondegenerate regular polytope of rank $r\geq9$ has at least
$2^{2r-1}$ flags, equivalently its string C-group has at least that order.
We use only this stated rank range and only connected string diagrams.
This is an external theorem, not a consequence of our finite diagnostics.

Apply Mar\'oti's original primitive-order theorem, including its lower
containment $A_m^\ell\leq G\leq S_m\wr S_\ell$ in the product-action
alternative, with $m\geq5$ and degree $n=\binom m k^\ell$.
If $\ell=1$, the natural $k=1$ action would give the excluded $A_n,S_n$.
For $k\geq2$, Whiston in the faithful $m$-point action gives
$d\leq m-1<n/2$. If $\ell\geq2$, the same abstract group has a faithful
transitive imprimitive action on $m\ell$ letters: the lower subgroup
$A_m^\ell$ is transitive in each block and primitivity of the product
action makes the coordinate action transitive. The induction hypothesis at the strictly smaller degree $m\ell<n$
provides the imprimitive bound and gives $d\leq m\ell/2+1<n/2$, since
$m\ell+2<m^\ell\leq\binom m k^\ell$ for $m\geq5,\ell\geq2$.

The Mathieu alternative is treated by the bounded checks below.
In the remaining alternative put
\[
 M(n)=n\prod_{i=0}^{\lfloor\log_2n\rfloor-1}(n-2^i).
\]
Exact integer arithmetic verifies $M(n)<2^{n-1}$ for $26\leq n\leq63$.
For $n\geq64$, with $q=\lfloor\log_2n\rfloor\geq6$, we have
$M(n)<n^{q+1}<2^{(q+1)^2}\leq2^{n-1}$; the inequality
$(q+1)^2\leq2^q-1$ holds at $q=6$ and persists by induction.
If $d\geq n/2$, then $d\geq13$, so Conder gives
$|G|\geq2^{2d-1}\geq2^{n-1}>M(n)$, a contradiction.
We use the exact product bound here, not the source's coarser logarithmic
estimate at its asserted small threshold.

\subsection{Complete finite screening}
For $2\leq n\leq25$, the required program visits every primitive
catalogue entry, omitting exactly $A_n,S_n$. It uses the following
necessary bounds, with all finite groups and actions computed exactly.
If $\Omega(t)$ counts prime factors with multiplicity, then
$d\leq\Omega(|G|)$, by the strictly increasing subgroup chain from IP.
All distinguished generators are involutions, so a proper subgroup
generated by all involutions excludes $G$ outright. For every possible
endpoint $x$, the subgroup on labels $2,\ldots,d-1$ commutes with $x$
and intersects $\gen x$ trivially. Consequently
\[
 d\leq1+\Omega(|C_G(x)|).
\]
Likewise the even labels generate an elementary abelian subgroup of
rank $\lceil d/2\rceil$. If every elementary abelian subgroup of a
Sylow two-subgroup has rank at most $e$, then $d\leq2e$.
The finite program computes this bound by the complete subgroup list
of that Sylow group; no assertion about its mere minimum number of
generators is substituted for its elementary abelian rank.

For primitive affine degree sixteen, the program verifies the regular
normal subgroup $2^4$. The normal-module lemma costs one generator,
and the faithful quotient embeds in $GL_4(2)\cong A_8$, whose independent
families have length at most six by Whiston. Thus $d\leq7<8$.
The exceptional isomorphism is explicitly checked and is also the
retained identification already used by common Topic B.
For the degree-nine affine group $AGL_2(3)$, the program verifies
$\Phi(GL_2(3))\cong C_2$ and the Frattini quotient $S_4$; its maximum
irredundant generating length is at most three. The irreducible
translation module costs one, giving $d\leq4<9/2$.
The degree-fifteen copy of $A_8$ has $d\leq6$ directly.

For both actions of $M_{11}$, the Sylow elementary abelian rank is two,
so $d\leq4$. The degree-twenty-three group $M_{23}$ is already excluded
by the endpoint-centralizer bound. For $M_{24}$ the two involution
centralizers have orders $7680,21504$. If $d\geq12$, the connected tail
on labels $2,\ldots,d-1$ has rank at least ten. Conder applied to that
tail, disjoint from $\gen x$, requires
$|C_G(x)|\geq2\cdot2^{2(d-2)-1}=2^{2d-4}\geq2^{20}$,
contradicting both actual centralizer orders.

The surviving actions are exactly the following GAP entries:
\[
 (6,1),(6,2),(7,5),(8,3),(8,5),(10,4),(10,5),(10,7),(12,2).
\]
The degree-eight copy $(8,4)$ is explicitly identified with $(7,5)$.
These are GAP's current indices, not Sims' indices printed in the source.
The maximum ranks searched are respectively $3,4,4,5,4,6,6,6,6$.
For $A_5,S_5$ they follow from Whiston; for $PSL_3(2)$ and $AGL_3(2)$
from the accepted common independent-length bounds; for the others
from the checked Sylow elementary abelian ranks two or three.

We enumerate connected involutory strings up to simultaneous normalizer
conjugation. At each prefix we discard dependence, failed string
commutations, or a failed endpoint-parabolic intersection. These are
necessary conditions for every subfamily of a C-marking. If $H$ is the
prefix subgroup and at least $t$ future generators remain after the next
one, their group centralizes $H$ and is disjoint from it; hence
$[C_G(H):Z(H)]\geq2^t$. This necessary inequality supplies a further
pruning rule. At a full generating tuple every pair of subset subgroups
is tested for IP. Stabilizer orbits at each extension preserve the full
prefix and cannot discard an ordered conjugacy class. No sample
markings, probabilistic generation test or source rank table is used.

The complete outputs, up to duality, are
\[
\begin{array}{c|c|l}
 n&G&\text{types at rank at least }n/2\\\hline
6&A_5&[3,5],\ [5,5]\\
6&S_5&[3,3,3],\ [4,5],\ [4,6],\ [5,6],\ [6,6]\\
10&S_6&[3,3,3,3].
\end{array}
\]
The last group is GAP $(10,5)$ and is identified exactly with $S_6$.
The full permutation tuples are retained, including both dual orientations.
Only the rank-four degree-six simplex exceeds half the permutation degree.
Thus this recovers the complete primitive branch of Theorem~1.2 and the
exception list of Proposition~3.2. The source's finer Table~3 maximum-rank
numbers are not imported as premises of this proof; any not independently
computed are separately attributed ancillary reports in the source map.
